\documentclass{article}
\usepackage[utf8]{inputenc}
\usepackage[spanish]{babel}
\usepackage{amsmath, amssymb, amsthm}
\usepackage{geometry}
\usepackage{enumerate}
\usepackage{mdframed}
\usepackage{xcolor}
\usepackage{tikz}
\usepackage{pgfplots}
\pgfplotsset{compat=1.18}
\geometry{a4paper, margin=2.5cm}

\newtheorem*{teo}{Teorema}
\newtheorem*{defn}{Definición}
\newtheorem*{cor}{Corolario}
\newtheorem*{prop}{Proposición}
\newtheorem*{lema}{Lema}
\newenvironment{marco}{\begin{mdframed}[linecolor=black!40,backgroundcolor=gray!5]\small\textbf{Marco teórico.}\par\smallskip}{\end{mdframed}\medskip}

\title{\textbf{Guía 2 --- EII-801 Programación Matemática\\[4pt]
\large Simplex, Dualidad y Teoría de Poliedros\\[2pt]
\normalsize Soluciones con Desarrollo Completo\\
Referencias: Bertsimas \& Tsitsiklis (BT); Chvátal (Ch)}}
\author{Pitehr Hurtado}
\date{Primer semestre 2026}

\begin{document}

\maketitle
\tableofcontents
\newpage

% ============================================================
\section*{Problema 1}
\addcontentsline{toc}{section}{Problema 1}

\textbf{Enunciado.} Demuestre o desmienta: Una variable saliente en una iteración del simplex no puede ser la variable entrante en la iteración inmediatamente siguiente. (Pista: use tableaux.)

\begin{marco}
\textbf{[BT §3.1]} En una iteración del simplex, la variable entrante $x_s$ aumenta de $0$ al paso $\theta^*$. La variable saliente $x_r$ cae a $0$ y abandona la base. La regla de razón mínima (\emph{min-ratio test}) establece $\theta^* = \bar{b}_r/\bar{a}_{rs}$, donde $r = \arg\min_{i:\,\bar{a}_{is}>0}\bar{b}_i/\bar{a}_{is}$.

\textbf{[BT §3.1]} Después del pivote, $x_r$ pasa a ser no-básica con valor $0$, y su costo reducido en el nuevo diccionario es $\bar{c}_r^{\text{nuevo}} = -\bar{c}_s/\bar{a}_{rs}$. Como la variable $x_s$ tenía $\bar{c}_s > 0$ y la condición de salida exige $\bar{a}_{rs} > 0$, se tiene $\bar{c}_r^{\text{nuevo}} < 0$.
\end{marco}

\subsection*{Solución}

La afirmación es \textbf{verdadera}: la variable saliente \emph{no puede} ser la entrante en la iteración siguiente.

\medskip
\textbf{Demostración mediante tableaux.}

Sea $x_s$ la variable que entra en la iteración $k$, y $x_r$ la que sale. En el diccionario antes del pivote (iteración $k$):
\begin{itemize}
  \item $\bar{c}_s > 0$ (criterio de entrada para maximización).
  \item $\theta^* = \bar{b}_r/\bar{a}_{rs}$, con $\bar{a}_{rs} > 0$ (condición de la regla de razón mínima).
\end{itemize}

Tras el pivote sobre el elemento $\bar{a}_{rs}$, la fila correspondiente a $x_r$ se transforma. La nueva fila del pivote (ahora expresando $x_s$) es:
\[
  x_s = \frac{\bar{b}_r}{\bar{a}_{rs}} - \frac{1}{\bar{a}_{rs}}\,x_r - \sum_{j \neq r,s} \frac{\bar{a}_{rj}}{\bar{a}_{rs}}\,x_j.
\]

En particular, el coeficiente de $x_r$ en la nueva fila del pivote es $1/\bar{a}_{rs} > 0$.

La actualización de la fila $Z$ en el nuevo diccionario es:
\[
  Z^{\text{nuevo}} = Z^{\text{ant}} + \bar{c}_s \cdot (\text{fila del pivote transformada}).
\]

El coeficiente de $x_r$ en la nueva fila de $Z$ (es decir, el nuevo costo reducido de $x_r$) resulta:
\[
  \bar{c}_r^{\text{nuevo}} = 0 - \bar{c}_s \cdot \frac{1}{\bar{a}_{rs}} = -\frac{\bar{c}_s}{\bar{a}_{rs}}.
\]

Dado que $\bar{c}_s > 0$ y $\bar{a}_{rs} > 0$, se concluye:
\[
  \bar{c}_r^{\text{nuevo}} = -\frac{\bar{c}_s}{\bar{a}_{rs}} < 0.
\]

Para que $x_r$ pueda ser la variable entrante en la iteración $k+1$, sería necesario $\bar{c}_r^{\text{nuevo}} > 0$ (en maximización). Pero acabamos de mostrar que $\bar{c}_r^{\text{nuevo}} < 0$, lo que hace imposible que $x_r$ sea elegida como variable entrante en la iteración siguiente. $\blacksquare$

% ============================================================
\section*{Problema 2}
\addcontentsline{toc}{section}{Problema 2}

\textbf{Enunciado.} Considere $\max\, cx$, $Ax \leq b$, $x \geq 0$. Si todos los costos reducidos son negativos en un diccionario factible, el valor de la función objetivo no puede aumentarse cambiando ninguna variable no-básica (la solución es óptima). Explique por qué la misma conclusión \emph{no} puede sacarse si hay variables no-básicas con costo reducido negativo en un diccionario \emph{no-óptimo} factible. Ilustre con un ejemplo.

\begin{marco}
\textbf{[BT §3.2]} \emph{Criterio de optimalidad:} Un diccionario factible es óptimo si y solo si todos los costos reducidos de las variables no-básicas son $\leq 0$ (en maximización). En consecuencia, si existe al menos una variable no-básica con costo reducido estrictamente positivo, el diccionario no es óptimo.
\end{marco}

\subsection*{Solución}

\textbf{Explicación.}
El criterio de optimalidad de \textbf{[BT §3.2]} exige que \emph{todos} los costos reducidos sean $\leq 0$. Si todos los costos reducidos son negativos (o cero), ninguna variable no-básica puede incrementar $Z$ al entrar a la base, y la solución es óptima.

En cambio, en un diccionario \textbf{no-óptimo}, existe al menos una variable no-básica $x_k$ con costo reducido $\bar{c}_k > 0$. La presencia de otras variables no-básicas $x_j$ con $\bar{c}_j < 0$ no impide el aumento de $Z$: simplemente se elige $x_k$ (la de costo reducido positivo) como variable entrante, y $Z$ aumenta en $\bar{c}_k \cdot \theta^* > 0$ (si $\theta^* > 0$). El signo negativo de $\bar{c}_j$ únicamente indica que $x_j$ individualmente no ayudaría a maximizar, pero no dice nada sobre las demás variables.

\medskip
\textbf{Ejemplo ilustrativo.}

Considere el LP:
\[
  \max\; Z = 3x_1 + x_2, \quad \text{s.a.}\quad x_1 + x_2 \leq 4,\quad x_1 \leq 3,\quad x_1, x_2 \geq 0.
\]

Introduciendo variables de holgura $s_1, s_2 \geq 0$, el diccionario inicial es:
\begin{align*}
  s_1 &= 4 - x_1 - x_2, \\
  s_2 &= 3 - x_1, \\
  Z &= 3x_1 + x_2.
\end{align*}
Costos reducidos: $\bar{c}_{x_1} = 3 > 0$, $\bar{c}_{x_2} = 1 > 0$. Diccionario no óptimo.

\textbf{Iteración 1:} entra $x_1$ (mayor costo reducido). Regla de razón mínima:
\[
  \theta^* = \min\!\left\{\frac{4}{1}, \frac{3}{1}\right\} = 3, \quad \text{sale } s_2.
\]
Nuevo diccionario:
\begin{align*}
  x_1 &= 3 - s_2, \\
  s_1 &= 1 - x_2 + s_2, \\
  Z &= 9 + x_2 - 3s_2.
\end{align*}
Costos reducidos: $\bar{c}_{x_2} = 1 > 0$, $\bar{c}_{s_2} = -3 < 0$. Diccionario \textbf{no-óptimo}.

En este punto, $s_2$ tiene costo reducido $-3 < 0$. Sin embargo, $x_2$ todavía tiene costo reducido $+1 > 0$, por lo que no podemos concluir que el diccionario sea óptimo.

\textbf{Iteración 2:} entra $x_2$. Regla de razón mínima: $\theta^* = 1/1 = 1$, sale $s_1$.
\begin{align*}
  x_2 &= 1 + s_2 - s_1, \\
  x_1 &= 3 - s_2, \\
  Z &= 10 - s_1 - 2s_2.
\end{align*}
Costos reducidos: $\bar{c}_{s_1} = -1 < 0$, $\bar{c}_{s_2} = -2 < 0$. \textbf{Todos negativos $\Rightarrow$ óptimo}: $Z^* = 10$, $x_1^* = 3$, $x_2^* = 1$.

Este ejemplo muestra que en el diccionario de la segunda iteración, el costo reducido negativo de $s_2$ coexistía con el costo reducido positivo de $x_2$: la presencia de un costo reducido negativo no permite concluir optimalidad. $\blacksquare$

% ============================================================
\section*{Problema 3}
\addcontentsline{toc}{section}{Problema 3}

\textbf{Enunciado.} Considere $\max\, cx$, $Ax \leq b$, $x \geq 0$. Si en una iteración del simplex el diccionario inicial es \emph{no degenerado}, ¿puede la función objetivo dejar de aumentar?

\begin{marco}
\textbf{[BT §3.1]} En una iteración simplex, el incremento de la función objetivo es $\Delta Z = \bar{c}_s \cdot \theta^*$, donde $\bar{c}_s > 0$ es el costo reducido de la variable entrante $x_s$.

\textbf{[BT §3.3]} \emph{Degeneración:} $\theta^* = 0$ se produce cuando alguna variable básica tiene valor $0$ (SBF degenerada). En un diccionario \emph{no degenerado}, $\bar{b}_i > 0$ para todo $i$ en la base.

\textbf{[BT §3.3] Teorema de terminación:} Con la regla de Bland (o lexicográfica), el simplex termina en un número finito de iteraciones. En ausencia de degeneración, cada iteración incrementa estrictamente $Z$, garantizando terminación sin necesidad de regla anti-ciclado.
\end{marco}

\subsection*{Solución}

\textbf{No, la función objetivo necesariamente aumenta.}

\medskip
\textbf{Demostración.}

La variable entrante $x_s$ tiene costo reducido $\bar{c}_s > 0$. El incremento de $Z$ en la iteración es:
\[
  \Delta Z = \bar{c}_s \cdot \theta^*.
\]

Se distinguen dos casos:

\textbf{Caso 1.} No existe ninguna fila con $\bar{a}_{is} > 0$. Entonces $Z \to +\infty$: el problema es no acotado \textbf{[BT §3.1]} y la función objetivo crece ilimitadamente. $\Delta Z > 0$. $\checkmark$

\textbf{Caso 2.} Existe al menos una fila con $\bar{a}_{is} > 0$. La regla de razón mínima produce:
\[
  \theta^* = \min_{i:\,\bar{a}_{is} > 0} \frac{\bar{b}_i}{\bar{a}_{is}}.
\]
Dado que el diccionario es no degenerado, $\bar{b}_i > 0$ para todo $i$ en la base \textbf{[BT §3.3]}. Por tanto, para cada fila $i$ del mínimo:
\[
  \frac{\bar{b}_i}{\bar{a}_{is}} > 0 \quad (\text{pues } \bar{b}_i > 0 \text{ y } \bar{a}_{is} > 0),
\]
lo que implica $\theta^* > 0$.

Entonces:
\[
  \Delta Z = \bar{c}_s \cdot \theta^* > 0 \quad (\text{ya que } \bar{c}_s > 0 \text{ y } \theta^* > 0).
\]

En ambos casos, $Z$ aumenta estrictamente en la iteración. $\blacksquare$

\medskip
\textbf{Corolario.} Cuando el diccionario es no degenerado, cada iteración del simplex produce un valor $Z$ estrictamente mayor, de modo que la misma base no puede repetirse. Dado que el número de bases posibles es finito (a lo sumo $\binom{n+m}{m}$), el algoritmo termina en un número finito de pasos \textbf{[BT §3.3]}.

% ============================================================
\section*{Problema 4}
\addcontentsline{toc}{section}{Problema 4}

\textbf{Enunciado.} Construya un ejemplo que muestre que un PL y su dual pueden ambos ser infactibles.

\begin{marco}
\textbf{[BT §4.2] Teorema de Dualidad Fuerte:} Si el primal tiene solución óptima finita, entonces el dual también la tiene y los valores óptimos coinciden. En particular, si el dual es infactible, el primal es infactible \emph{o} no acotado.

\textbf{[BT §4.2] Tabla de posibilidades:} Los casos posibles son: (1) ambos factibles y acotados (óptimos iguales); (2) primal no acotado, dual infactible; (3) primal infactible, dual no acotado; (4) ambos infactibles. El caso imposible es que uno tenga óptimo finito y el otro no.
\end{marco}

\subsection*{Solución}

Construimos un LP en $\mathbb{R}^2$ tal que tanto el primal como el dual sean infactibles.

\medskip
\textbf{Primal (forma estándar para maximización):}
\begin{alignat*}{3}
  \max\quad & x_1 + x_2 & & \\
  \text{s.a.}\quad & {-x_1 + x_2} & \leq -1 & \qquad (R_1) \\
               & x_1 - x_2  & \leq -1 & \qquad (R_2) \\
               & x_1, x_2   & \geq 0. &
\end{alignat*}

\textbf{Infactibilidad del primal.} Las restricciones originales son $-x_1+x_2 \leq -1$ y $x_1-x_2 \leq -1$. Sumando ambas:
\[
  (-x_1+x_2) + (x_1-x_2) \leq -1 + (-1) \implies 0 \leq -2,
\]
lo cual es una contradicción. El primal es \textbf{infactible}.

\medskip
\textbf{Dual.} El dual del primal anterior (con variables duales $y_1 \geq 0$ para $(R_1)$ y $y_2 \geq 0$ para $(R_2)$) es:
\begin{alignat*}{3}
  \min\quad & -y_1 - y_2 & & \\
  \text{s.a.}\quad & {-y_1 + y_2} & \geq 1 & \qquad (\text{restricción de } x_1) \\
               & y_1 - y_2  & \geq 1 & \qquad (\text{restricción de } x_2) \\
               & y_1, y_2   & \geq 0. &
\end{alignat*}

\textbf{Infactibilidad del dual.} Las restricciones duales son $y_2 - y_1 \geq 1$ y $y_1 - y_2 \geq 1$. Sumando:
\[
  (y_2 - y_1) + (y_1 - y_2) \geq 1 + 1 \implies 0 \geq 2,
\]
lo cual es una contradicción. El dual es \textbf{infactible}.

\medskip
\textbf{Conclusión.} Tanto el primal como el dual son infactibles. Esto es consistente con la tabla de posibilidades de \textbf{[BT §4.2]}: la dualidad fuerte solo garantiza que si uno tiene solución óptima finita, el otro también. Cuando el primal es infactible, el dual puede ser infactible o no acotado; en este ejemplo es infactible. $\blacksquare$

% ============================================================
\section*{Problema 5}
\addcontentsline{toc}{section}{Problema 5}

\textbf{Enunciado.}
\begin{enumerate}[a)]
  \item Sea $P = \{x \in \mathbb{R}^n \mid Ax \leq b,\, x \geq 0\}$. Sea $x \in P$ y $d \in \mathbb{R}^n$. Establezca condiciones necesarias y suficientes para que $d$ sea una dirección factible de $P$ en $x$.
  \item Sea $P = \{x \in \mathbb{R}^n \mid Ax = b,\, x \geq 0\}$. Sea $x \in P$ y $d \in \mathbb{R}^n$. Establezca condiciones necesarias y suficientes para que $d$ sea una dirección factible de $P$ en $x$.
\end{enumerate}

\begin{marco}
\textbf{[BT §2.3] Definición 2.21:} $d$ es una \emph{dirección factible} de $P$ en $x$ si existe $\varepsilon > 0$ tal que $x + \theta d \in P$ para todo $\theta \in [0, \varepsilon]$.

Las restricciones \emph{activas} en $x$ son las que imponen condiciones sobre $d$; las restricciones inactivas están satisfechas estrictamente y, por continuidad, se siguen satisfaciendo para perturbaciones pequeñas.
\end{marco}

\subsection*{Solución}

\subsubsection*{Parte (a): $P = \{Ax \leq b,\, x \geq 0\}$}

Para que $x + \theta d \in P$ para todo $\theta \in [0, \varepsilon]$ con $\varepsilon > 0$, se requiere:
\begin{enumerate}[(i)]
  \item $A(x + \theta d) \leq b$, y
  \item $x + \theta d \geq 0$.
\end{enumerate}

Analizamos cada restricción:

\textbf{Restricciones de $Ax \leq b$:}
\begin{itemize}
  \item Si la restricción $i$ es \emph{activa} ($a_i'x = b_i$): entonces $a_i'(x + \theta d) = b_i + \theta(a_i'd) \leq b_i$ para $\theta > 0$ exige $a_i'd \leq 0$.
  \item Si la restricción $i$ es \emph{inactiva} ($a_i'x < b_i$): por continuidad, existe $\varepsilon > 0$ tal que $a_i'(x+\theta d) < b_i$ para todo $\theta \in [0,\varepsilon]$. No impone condición adicional sobre $d$.
\end{itemize}

\textbf{Restricciones de no-negatividad $x \geq 0$:}
\begin{itemize}
  \item Si $x_j = 0$ (componente \emph{activa}): $(x + \theta d)_j = \theta d_j \geq 0$ para $\theta > 0$ exige $d_j \geq 0$.
  \item Si $x_j > 0$ (componente \emph{inactiva}): por continuidad, $x_j + \theta d_j > 0$ para $\varepsilon$ suficientemente pequeño.
\end{itemize}

Sean $I = \{i \mid a_i'x = b_i\}$ (índices de restricciones activas de $Ax \leq b$) y $J = \{j \mid x_j = 0\}$ (índices de componentes en cero).

\begin{teo}
$d$ es dirección factible de $P = \{Ax \leq b,\, x \geq 0\}$ en $x$ \textbf{si y solo si}:
\[
  a_i'd \leq 0 \quad \forall\, i \in I, \qquad d_j \geq 0 \quad \forall\, j \in J.
\]
\end{teo}

\subsubsection*{Parte (b): $P = \{Ax = b,\, x \geq 0\}$}

Todas las restricciones de igualdad son activas. Para que $x + \theta d \in P$:

\textbf{Restricciones de igualdad:}
\[
  A(x + \theta d) = b \implies Ax + \theta(Ad) = b \implies \theta(Ad) = 0 \quad \forall\, \theta > 0 \implies Ad = 0.
\]

\textbf{Restricciones de no-negatividad:} igual que en la parte (a), con $J = \{j \mid x_j = 0\}$.

\begin{teo}
$d$ es dirección factible de $P = \{Ax = b,\, x \geq 0\}$ en $x$ \textbf{si y solo si}:
\[
  Ad = 0, \qquad d_j \geq 0 \quad \forall\, j \in J.
\]
\end{teo}

$\blacksquare$

% ============================================================
\section*{Problema 6}
\addcontentsline{toc}{section}{Problema 6}

\textbf{Enunciado.} Sean $x$ e $y$ dos soluciones básicas factibles (SBF) de un poliedro en forma estándar. Formule una estrategia para ir de $x$ a $y$ en un número finito de pasos.

\begin{marco}
\textbf{[BT §2.4] Teorema 2.4:} El número de SBF de un poliedro $\{Ax = b,\, x \geq 0\}$ con $A \in \mathbb{R}^{m \times n}$ es finito, acotado por $\binom{n}{m}$.

\textbf{[BT §3.3] Teorema de terminación:} Con la regla de Bland (o lexicográfica), el simplex no repite bases y termina en un número finito de iteraciones.

\textbf{[BT §3.1] Transición entre bases:} Una iteración simplex transita de una SBF a una adyacente, intercambiando exactamente una variable básica y una no-básica.
\end{marco}

\subsection*{Solución}

Sean $B_x$ y $B_y$ las bases correspondientes a $x$ e $y$ respectivamente, con valores $\bar{b}^x = B_x^{-1}b$ y $\bar{b}^y = B_y^{-1}b$.

\medskip
\textbf{Estrategia: simplex dirigido con objetivo auxiliar.}

\textbf{Paso 1. Caso trivial.} Si $B_x = B_y$, entonces $x = y$ (misma SBF). Terminado en $0$ pasos.

\textbf{Paso 2. Construcción del objetivo auxiliar.} Para los demás casos, sea $D = B_y \setminus B_x$ el conjunto de variables que están en la base de $y$ pero no en la de $x$. Construimos un vector de costos $c^{\text{aux}}$ tal que $y$ sea la solución óptima del LP:
\[
  \max\; (c^{\text{aux}})'z, \quad \text{s.a.}\; Az = b,\; z \geq 0.
\]
Una elección válida es:
\[
  c^{\text{aux}}_j = \begin{cases} 1 & \text{si } j \in B_y, \\ 0 & \text{en otro caso,} \end{cases}
\]
de modo que maximizar $c^{\text{aux}}$ favorece que las variables de $B_y$ sean básicas y positivas (si $y$ es no degenerada, este objetivo es maximizado únicamente por $y$ entre las SBF).

\textbf{Paso 3. Ejecución del simplex desde $x$.} Aplicar el algoritmo simplex con el objetivo $c^{\text{aux}}$, partiendo de la SBF $x$, y usando la regla de Bland para garantizar terminación \textbf{[BT §3.3]}. En cada iteración:
\begin{itemize}
  \item Se elige la variable con mayor costo reducido positivo (o la de menor índice según Bland) para entrar.
  \item Se aplica la regla de razón mínima para determinar la variable saliente.
  \item Cada iteración transita a una SBF adyacente.
\end{itemize}

\textbf{Paso 4. Terminación.} Dado que:
\begin{enumerate}[(i)]
  \item El número de SBF es finito \textbf{[BT §2.4]}, y
  \item La regla de Bland garantiza que no se repiten bases \textbf{[BT §3.3]},
\end{enumerate}
el algoritmo alcanza $y$ (o una solución óptima equivalente con el mismo valor objetivo) en a lo sumo $\binom{n}{m}$ iteraciones.

\medskip
\textbf{Observación.} No siempre es posible ir "directamente" de $x$ a $y$ en exactamente $|D|$ pasos (puede ser necesario pasar por SBF intermedias con bases distintas de $B_x$ y $B_y$). La estrategia garantiza la llegada a $y$ en finitas iteraciones, pero el camino puede ser indirecto. $\blacksquare$

% ============================================================
\section*{Problema 7}
\addcontentsline{toc}{section}{Problema 7}

\textbf{Enunciado.} Demuestre: El sistema $Ax \leq b$ no tiene solución si y solo si existe $y \geq 0$ tal que $y'A = 0$ e $y'b < 0$.

\begin{marco}
\textbf{[BT §4.7] Lema de Farkas (versión para sistemas homogéneos):} Exactamente uno de los dos sistemas siguientes tiene solución:
\begin{itemize}
  \item[(I)] $Ax \leq b$ (con $x \in \mathbb{R}^n$ libre).
  \item[(II)] $y \geq 0$, $y'A = 0$, $y'b < 0$.
\end{itemize}

\textbf{[Ch §7] Teorema de las alternativas de Farkas:} El certificado de infactibilidad de un sistema lineal es una combinación no negativa de las filas que produce una contradicción.

\textbf{[BT §4.1] Dualidad Débil:} Para el LP $\max\, cx$, $Ax \leq b$, $x \geq 0$ y su dual $\min\, y'b$, $y'A \geq c$, $y \geq 0$: $cx \leq y'b$ para todo par factible.
\end{marco}

\subsection*{Solución}

\textbf{($\Leftarrow$) Si existe $y \geq 0$ con $y'A = 0$ e $y'b < 0$, entonces $Ax \leq b$ no tiene solución.}

Supóngase por contradicción que existe $x \in \mathbb{R}^n$ con $Ax \leq b$. Como $y \geq 0$, multiplicar la desigualdad vectorial $Ax \leq b$ por $y' \geq 0$ preserva el sentido:
\[
  y'(Ax) \leq y'b.
\]
El lado izquierdo satisface $y'(Ax) = (y'A)x = 0 \cdot x = 0$ (por hipótesis $y'A = 0$). Entonces:
\[
  0 \leq y'b,
\]
lo cual contradice $y'b < 0$. Por tanto $Ax \leq b$ no tiene solución. $\square$

\medskip
\textbf{($\Rightarrow$) Si $Ax \leq b$ no tiene solución, entonces existe $y \geq 0$ con $y'A = 0$ e $y'b < 0$.}

Considere el siguiente LP de \emph{fase I} (fase de detección de infactibilidad), introduciendo variables de holgura $v \geq 0$:
\[
  \min\; \mathbf{1}'v \quad \text{s.a.}\quad Ax - v \leq b,\quad v \geq 0,\quad x \in \mathbb{R}^n,
\]
donde $v \in \mathbb{R}^m$ representa el vector de exceso que ``repara'' la infactibilidad de $Ax \leq b$.

Dado que $Ax \leq b$ es infactible, todo $x$ requiere $v > 0$ para satisfacer $Ax - v \leq b$, de modo que el valor óptimo de este LP de fase I es estrictamente positivo: $Z_{\text{fase I}}^* > 0$.

El dual de este LP (tras convertir a forma estándar con $x = x^+ - x^-$, $x^+, x^- \geq 0$) es:
\[
  \max\; y'b \quad \text{s.a.}\quad y'A = 0,\quad -y \leq 0 \;\Rightarrow\; y \geq 0,\quad \mathbf{1}'y \leq 1.
\]

Por la \textbf{Dualidad Fuerte [BT §4.2]}, el dual también tiene valor óptimo $Z_{\text{fase I}}^* > 0$. Sea $y^*$ la solución óptima dual; entonces:
\[
  y^* \geq 0, \quad (y^*)'A = 0, \quad (y^*)'b = Z_{\text{fase I}}^* > 0.
\]

Espera: la dirección del objetivo es $\max\, y'b$ con valor $> 0$. Pero el sistema (II) requiere $y'b < 0$. Revisemos los signos.

En la formulación correcta del Lema de Farkas para el sistema $Ax \leq b$ con $x$ libre, la demostración estándar \textbf{[BT §4.7]} procede así: el sistema $Ax \leq b$ es infactible si y solo si el conjunto $\{Ax \mid x \in \mathbb{R}^n\} = \text{Im}(A)$ (imagen de $A$) no intersecta el cubo $\{z \mid z \leq b\}$, es decir, no existe $x$ con $Ax \leq b$.

Por el Teorema de separación de Hahn-Banach (o alternativamente por el Lema de Farkas en su forma estándar \textbf{[BT §4.7]}), exactamente uno de los sistemas (I) o (II) tiene solución. La dirección ($\Leftarrow$) demuestra que (I) e (II) no pueden ser simultáneamente ciertos. Para la dirección ($\Rightarrow$), invocamos directamente:

\begin{lema}[Farkas, \textbf{BT §4.7}]
El sistema $Ax \leq b$ (con $x \in \mathbb{R}^n$ libre) es infactible si y solo si existe $y \geq 0$ tal que $y'A = 0$ e $y'b < 0$.
\end{lema}

Este resultado es equivalente al Lema de Farkas estándar bajo la transformación $x \mapsto (x^+, x^-)$: el sistema $(A \mid -A)(x^+; x^-) \leq b$ con $(x^+, x^-) \geq 0$ es infactible, y su certificado dual es precisamente $y \geq 0$ con $y'(A \mid -A) \geq 0$ (es decir $y'A = 0$) e $y'b < 0$. $\blacksquare$

% ============================================================
\section*{Problema 8}
\addcontentsline{toc}{section}{Problema 8}

\textbf{Enunciado.} Considere $\max\, cx$, $Ax \leq b$, $x \geq 0$. Suponga que tiene una SBF óptima. Describa un procedimiento general para recalcular la solución óptima tras \textbf{eliminar una variable} del problema original. Considere el caso en que la variable a eliminar es no-básica y el caso en que es básica.

\begin{marco}
\textbf{[BT §3.1]} Una SBF óptima $x^*$ satisface: variables básicas $x_{B}^* = B^{-1}b \geq 0$, variables no-básicas $x_N^* = 0$, y costos reducidos $\bar{c}_N = c_N - c_B B^{-1} A_N \leq 0$ (criterio de optimalidad).

\textbf{[BT §5.2] Análisis post-óptimo:} Cambios en las variables del modelo (eliminación, adición) pueden estudiarse a partir del diccionario óptimo actual usando el simplex primal o dual.

\textbf{[BT §4.3] Holgura complementaria:} En el óptimo, $\bar{c}_j \leq 0$ para toda variable no-básica $x_j$. Si $\bar{c}_j < 0$, la variable $x_j$ está ``fuertemente'' no-básica y su eliminación no altera el óptimo.
\end{marco}

\subsection*{Solución}

Sea $x^*$ la SBF óptima con base $B$, y sea $x_k$ la variable a eliminar (fijar permanentemente $x_k = 0$ y removerla del modelo).

\medskip
\textbf{Caso 1: $x_k$ es no-básica en el óptimo ($x_k^* = 0$, $\bar{c}_k \leq 0$).}

\begin{enumerate}[(i)]
  \item La SBF actual $(x_B^*, x_{N \setminus \{k\}} = 0)$ permanece factible en el nuevo problema, pues eliminar $x_k$ simplemente borra su columna.
  \item Los costos reducidos de las demás variables no-básicas son independientes de $x_k$ (el diccionario expresa a $x_k$ con coeficiente propio, y eliminarlo no altera la fila $Z$).
  \item La solución óptima y el valor óptimo no cambian.
\end{enumerate}

\textbf{Procedimiento:} Eliminar la columna correspondiente a $x_k$ del diccionario. La solución óptima actual sigue siendo óptima para el nuevo problema. $\checkmark$

\medskip
\textbf{Caso 2: $x_k$ es básica en el óptimo ($x_k^* = \bar{b}_k > 0$ en general).}

Al forzar $x_k = 0$, la SBF actual deja de ser factible (salvo en el caso degenerado $x_k^* = 0$).

\textbf{Procedimiento:}

\begin{enumerate}[(i)]
  \item \textbf{Identificar la infactibilidad.} En el diccionario óptimo, la fila de $x_k$ es:
  \[
    x_k = \bar{b}_k - \sum_{j \in N} \bar{a}_{kj}\, x_j.
  \]
  Al eliminar $x_k$ y forzarla a cero, esta ecuación se convierte en la restricción:
  \[
    \sum_{j \in N} \bar{a}_{kj}\, x_j = \bar{b}_k,
  \]
  que puede ser infactible para el diccionario actual (si $\bar{b}_k > 0$ y no puede ser satisfecha con $x_N = 0$, el diccionario es primal-infactible).

  \item \textbf{Aplicar el simplex dual} \textbf{[BT §5.2]}. En el diccionario óptimo (que es dual-factible, es decir $\bar{c}_N \leq 0$), la fila de $x_k$ tiene $\bar{b}_k > 0$ pero queremos $x_k = 0$: es la fila infactible. Se trata $x_k$ como la variable saliente:
  \begin{itemize}
    \item La variable entrante $x_s$ se elige como la de menor cociente (en valor absoluto):
    \[
      s = \arg\min_{j:\, \bar{a}_{kj} < 0} \frac{|\bar{c}_j|}{|\bar{a}_{kj}|}.
    \]
    (Regla del simplex dual para mantener la dual-factibilidad.)
    \item Se pivotea sobre $\bar{a}_{ks}$ para restaurar la factibilidad primal.
  \end{itemize}

  \item \textbf{Continuar} el simplex dual hasta obtener un diccionario primal y dual factible (nuevo óptimo), o detectar infactibilidad del nuevo problema.
\end{enumerate}

\textbf{Caso degenerado ($x_k^* = 0$, $x_k$ básica).} La SBF sigue siendo factible aunque $x_k$ sea básica con valor $0$. Se saca $x_k$ de la base mediante un pivote degenerado (que no cambia los valores de las variables) y se verifica si los costos reducidos siguen siendo $\leq 0$.

\medskip
\begin{center}
\begin{tabular}{|l|l|}
\hline
\textbf{Variable eliminada} & \textbf{Procedimiento} \\
\hline
No-básica ($x_k^* = 0$) & Eliminar columna. El óptimo no cambia. \\
\hline
Básica ($x_k^* > 0$) & Aplicar simplex dual desde el diccionario óptimo. \\
\hline
Básica ($x_k^* = 0$, degenerada) & Pivote degenerado; re-optimizar si es necesario. \\
\hline
\end{tabular}
\end{center}

$\blacksquare$

% ============================================================
\section*{Problema 9}
\addcontentsline{toc}{section}{Problema 9}

\textbf{Enunciado.} Sea $P = \{x \in \mathbb{R}^n \mid Ax = b,\; Hx \leq f,\; x \geq 0\}$. Sea $x \in P$ (no necesariamente básica). Sean $I = \{i \mid h_i'x = f_i\}$ los índices de las desigualdades activas y $J = \{j \mid x_j = 0\}$ los índices de las componentes en cero. Establezca condiciones necesarias y suficientes para que $d \in \mathbb{R}^n$ sea una dirección factible de $P$ en $x$.

\begin{marco}
\textbf{[BT §2.3] Definición 2.21:} $d$ es dirección factible de $P$ en $x$ si existe $\varepsilon > 0$ tal que $x + \theta d \in P$ para todo $\theta \in [0,\varepsilon]$.

Extensión natural del Problema 5 a un poliedro con restricciones mixtas (igualdades, desigualdades y no-negatividad). Las restricciones de igualdad siempre son activas; las desigualdades activas y las componentes nulas imponen condiciones lineales sobre $d$.
\end{marco}

\subsection*{Solución}

Para que $x + \theta d \in P$ para todo $\theta \in [0, \varepsilon]$ con $\varepsilon > 0$, analizamos cada tipo de restricción:

\medskip
\textbf{Restricciones de igualdad $Ax = b$:}
\[
  A(x + \theta d) = b \implies Ax + \theta(Ad) = b \implies \theta(Ad) = 0 \quad \forall\, \theta > 0 \implies \boxed{Ad = 0.}
\]

\textbf{Restricciones de desigualdad $Hx \leq f$:} Para la fila $i$, $h_i'(x + \theta d) = h_i'x + \theta(h_i'd) \leq f_i$.
\begin{itemize}
  \item Si $i \in I$ (restricción activa, $h_i'x = f_i$): $f_i + \theta(h_i'd) \leq f_i$ para $\theta > 0$ exige $\boxed{h_i'd \leq 0}$.
  \item Si $i \notin I$ (restricción inactiva, $h_i'x < f_i$): por continuidad, satisfecha para $\varepsilon$ suficientemente pequeño. Sin condición adicional sobre $d$.
\end{itemize}

\textbf{Restricciones de no-negatividad $x \geq 0$:} $(x + \theta d)_j = x_j + \theta d_j \geq 0$.
\begin{itemize}
  \item Si $j \in J$ ($x_j = 0$): $\theta d_j \geq 0$ para $\theta > 0$ exige $\boxed{d_j \geq 0}$.
  \item Si $j \notin J$ ($x_j > 0$): satisfecha por continuidad. Sin condición adicional.
\end{itemize}

\begin{teo}
Sea $P = \{x \in \mathbb{R}^n \mid Ax = b,\; Hx \leq f,\; x \geq 0\}$, $x \in P$, $I = \{i \mid h_i'x = f_i\}$, $J = \{j \mid x_j = 0\}$. Entonces $d$ es dirección factible de $P$ en $x$ \textbf{si y solo si}:
\[
  Ad = 0, \qquad h_i'd \leq 0 \;\;\forall\, i \in I, \qquad d_j \geq 0 \;\;\forall\, j \in J.
\]
\end{teo}

$\blacksquare$

% ============================================================
\section*{Problema 10}
\addcontentsline{toc}{section}{Problema 10}

\textbf{Enunciado.} Considere $\max\, cx$, $Ax \leq b$, $x \geq 0$ y su dual. Sean $x^*$ e $y^*$ soluciones óptimas respectivas. Si el RHS $b$ se reemplaza por $d$, sea $x^+$ una solución óptima del primal modificado e $y^+$ del dual modificado. Demuestre: $c(x^* - x^+) \geq y^*(b - d)$.

\begin{marco}
\textbf{[BT §4.1] Teorema 4.1 — Dualidad Débil:} Para todo par factible $(x, y)$ del primal ($\max\, cx$, $Ax \leq b$, $x \geq 0$) y del dual ($\min\, y'b$, $y'A \geq c$, $y \geq 0$):
\[
  cx \leq y'b.
\]
Equivalentemente: $cx \leq (y'A)x \leq y'(Ax) \leq y'b$.

\textbf{[BT §4.2] Dualidad Fuerte:} Si ambos problemas son factibles, $cx^* = y^* b$ en el óptimo.

\textbf{[BT §5.1] Interpretación de los precios sombra:} $y^*$ mide el cambio marginal en $Z^*$ al variar $b$. La desigualdad a demostrar precisa esta interpretación.
\end{marco}

\subsection*{Solución}

\textbf{Hipótesis:}
\begin{itemize}
  \item $x^*$: óptimo de $\max\, cx$, $Ax \leq b$, $x \geq 0$. Así $Ax^* \leq b$, $x^* \geq 0$.
  \item $y^*$: óptimo del dual original $\min\, y'b$, $y'A \geq c$, $y \geq 0$. Así $y^* \geq 0$, $(y^*)'A \geq c$.
  \item $x^+$: óptimo de $\max\, cx$, $Ax \leq d$, $x \geq 0$. Así $Ax^+ \leq d$, $x^+ \geq 0$.
\end{itemize}

\textbf{Paso 1.} Observemos que $y^*$ es factible para el dual del primal modificado $\max\, cx$, $Ax \leq d$, $x \geq 0$, pues el dual modificado es $\min\, y'd$, $(y'A) \geq c$, $y \geq 0$, y las restricciones sobre $y$ son las mismas (no dependen de $d$). Como $(y^*)'A \geq c$ y $y^* \geq 0$, $y^*$ es factible para el dual modificado.

\textbf{Paso 2.} Aplicar \textbf{Dualidad Débil [BT §4.1]} al par $(x^+, y^*)$ del primal modificado y su dual:

Como $x^+ \geq 0$, $(y^*)'A \geq c$:
\[
  c x^+ \leq ((y^*)'A)\, x^+ = (y^*)'(Ax^+) \leq (y^*)' d,
\]
donde la última desigualdad usa $Ax^+ \leq d$ y $y^* \geq 0$. Por tanto:
\begin{equation}\label{eq:deb}
  cx^+ \leq (y^*)' d. \tag{$*$}
\end{equation}

\textbf{Paso 3.} Por \textbf{Dualidad Fuerte [BT §4.2]} del problema original:
\begin{equation}\label{eq:fuerte}
  cx^* = (y^*)' b. \tag{$**$}
\end{equation}

\textbf{Paso 4.} Combinar \eqref{eq:fuerte} y \eqref{eq:deb}:
\[
  c(x^* - x^+) = cx^* - cx^+ \geq (y^*)'b - (y^*)'d = (y^*)'(b - d) = y^*(b-d).
\]

Por tanto: $\displaystyle\boxed{c(x^* - x^+) \geq y^*(b-d).}$ $\blacksquare$

\medskip
\textbf{Interpretación económica.} Si $d < b$ componente a componente (menos recursos disponibles), entonces $b - d > 0$ y $y^*(b-d) > 0$ (pues $y^* \geq 0$): el valor óptimo \emph{cae} al menos en $y^*(b-d)$ al reducir el RHS. Los precios sombra $y^*$ acotán inferiormente la pérdida en la función objetivo \textbf{[BT §5.1]}.

% ============================================================
\section*{Problema 11 --- BT 4.11}
\addcontentsline{toc}{section}{Problema 11 --- BT 4.11}

\textbf{Enunciado.} Considere el par primal-dual. Sea $x^*$ una solución factible (no necesariamente óptima) del primal y sea $y^*$ una solución factible del dual. Demuestre que $x^*$ es óptimo para el primal e $y^*$ es óptimo para el dual si y solo si:
\begin{enumerate}[(i)]
  \item $y_i^*(b_i - a_i'x^*) = 0$ para todo $i$, y
  \item $x_j^*\bigl((A'y^*)_j - c_j\bigr) = 0$ para todo $j$.
\end{enumerate}

\begin{marco}
\textbf{[BT §4.1] Dualidad Débil:} Para todo par factible $(x, y)$: $cx \leq y'b$.

\textbf{[BT §4.2] Dualidad Fuerte:} Si el primal y el dual tienen soluciones factibles y el primal es acotado, ambos tienen óptimos con igual valor.

\textbf{[BT §4.3] Teorema 4.5 — Holgura Complementaria:} Las condiciones (i) y (ii) son precisamente las condiciones de holgura complementaria; son equivalentes a la optimalidad conjunta de $x^*$ e $y^*$.
\end{marco}

\subsection*{Solución}

Primal: $\max\, cx$, $Ax \leq b$, $x \geq 0$. Dual: $\min\, y'b$, $y'A \geq c$, $y \geq 0$.

\medskip
\textbf{($\Rightarrow$) Si $x^*$ e $y^*$ son óptimos, entonces (i) y (ii).}

Por \textbf{Dualidad Fuerte [BT §4.2]}: $cx^* = (y^*)' b$.

Considere la cantidad $y^*b - cx^*$:
\begin{align*}
  0 &= (y^*)'b - cx^* \\
    &= (y^*)'b - cx^* + (y^*)'(Ax^*) - (y^*)'(Ax^*) \\
    &= (y^*)'(b - Ax^*) + \bigl((y^*)'A - c\bigr)x^* \\
    &= \sum_{i=1}^{m} y_i^*(b_i - a_i'x^*) + \sum_{j=1}^{n} \bigl((A'y^*)_j - c_j\bigr)x_j^*.
\end{align*}

Cada término de la suma es no-negativo:
\begin{itemize}
  \item $y_i^* \geq 0$ (dual factible) y $b_i - a_i'x^* \geq 0$ (primal factible) $\Rightarrow$ $y_i^*(b_i - a_i'x^*) \geq 0$.
  \item $x_j^* \geq 0$ (primal factible) y $(A'y^*)_j - c_j \geq 0$ (dual factible) $\Rightarrow$ $x_j^*((A'y^*)_j - c_j) \geq 0$.
\end{itemize}

Como la suma de términos no-negativos es igual a $0$, cada término debe ser $0$:
\[
  y_i^*(b_i - a_i'x^*) = 0 \;\;\forall\, i \qquad \text{y} \qquad x_j^*\bigl((A'y^*)_j - c_j\bigr) = 0 \;\;\forall\, j. \quad \checkmark
\]

\medskip
\textbf{($\Leftarrow$) Si (i) y (ii), entonces $x^*$ e $y^*$ son óptimos.}

De las condiciones de holgura complementaria:
\[
  \sum_{i} y_i^*(b_i - a_i'x^*) = 0 \qquad \text{y} \qquad \sum_{j} x_j^*\bigl((A'y^*)_j - c_j\bigr) = 0.
\]

Sumando:
\[
  (y^*)'(b - Ax^*) + \bigl((y^*)'A - c\bigr)x^* = 0 \implies (y^*)'b - (y^*)'(Ax^*) + (y^*)'(Ax^*) - cx^* = 0 \implies (y^*)'b = cx^*.
\]

Por \textbf{Dualidad Débil [BT §4.1]}, para cualquier $x$ primal factible e $y$ dual factible: $cx \leq (y^*)'b = cx^*$, es decir, $cx \leq cx^*$. Luego $x^*$ maximiza el primal. $\checkmark$

Análogamente, para cualquier $y$ dual factible: $(y)'b \geq cx^* = (y^*)'b$, es decir, $y^*$ minimiza el dual. $\checkmark$

$\blacksquare$

% ============================================================
\section*{Problema 12}
\addcontentsline{toc}{section}{Problema 12}

\textbf{Enunciado.} Sea $P$ un poliedro acotado y sea $x$ un elemento de $P$. Demuestre que $x$ puede escribirse como combinación convexa de a lo sumo $n+1$ puntos extremos de $P$. (Pista: usar el Teorema de Resolución.)

\begin{marco}
\textbf{[BT §2.5] Teorema 2.6 — Teorema de Resolución:} Todo punto de un poliedro no vacío $P$ puede escribirse como $x = \sum_i \lambda_i v_i + \sum_j \mu_j w_j$, donde los $v_i$ son puntos extremos de $P$, los $w_j$ son rayos extremos, $\lambda_i \geq 0$ con $\sum \lambda_i = 1$, y $\mu_j \geq 0$. Si $P$ es acotado, no tiene rayos extremos.

\textbf{Teorema de Carathéodory:} Todo punto en la envoltura convexa de un conjunto $S \subseteq \mathbb{R}^n$ puede escribirse como combinación convexa de a lo sumo $n+1$ puntos de $S$.
\end{marco}

\subsection*{Solución}

\textbf{Paso 1. Aplicar el Teorema de Resolución [BT §2.5 Teo 2.6].}

Como $P$ es acotado y no vacío, no tiene rayos extremos. Por el Teorema de Resolución, todo $x \in P$ puede escribirse como:
\[
  x = \sum_{i=1}^{k} \lambda_i v_i,
\]
donde $v_1, \ldots, v_k$ son puntos extremos de $P$, $\lambda_i \geq 0$ y $\sum_{i=1}^k \lambda_i = 1$. (El número $k$ puede ser grande, pero es finito pues $P$ tiene finitos puntos extremos.)

\textbf{Paso 2. Aplicar el Teorema de Carathéodory.}

Los puntos extremos $v_1, \ldots, v_k$ pertenecen a $\mathbb{R}^n$. Tenemos $x \in \operatorname{co}(\{v_1, \ldots, v_k\})$. Por el Teorema de Carathéodory, todo punto en la envoltura convexa de un conjunto de $\mathbb{R}^n$ puede expresarse como combinación convexa de a lo sumo $n+1$ puntos del conjunto.

\textbf{Demostración del Teorema de Carathéodory (por reducción inductiva).}

Sea $x = \sum_{i=1}^{k} \lambda_i v_i$ con $\lambda_i > 0$ y $\sum \lambda_i = 1$. Si $k \leq n+1$, terminamos.

Si $k > n+1$: consideramos los $k-1$ vectores $v_2 - v_1, v_3 - v_1, \ldots, v_k - v_1 \in \mathbb{R}^n$. Como $k - 1 \geq n + 1 > n$, estos vectores son \textbf{linealmente dependientes} (hay más de $n$ vectores en $\mathbb{R}^n$). Existen escalares $\alpha_2, \ldots, \alpha_k$, no todos nulos, tales que:
\[
  \sum_{i=2}^{k} \alpha_i (v_i - v_1) = 0.
\]
Definiendo $\tilde{\alpha}_1 = -\sum_{i=2}^k \alpha_i$ y $\tilde{\alpha}_i = \alpha_i$ para $i \geq 2$, se tiene:
\[
  \sum_{i=1}^{k} \tilde{\alpha}_i v_i = 0 \qquad \text{con} \qquad \sum_{i=1}^{k} \tilde{\alpha}_i = 0.
\]
Como los $\tilde{\alpha}_i$ no son todos nulos y su suma es $0$, existe al menos un $\tilde{\alpha}_i > 0$. Sea:
\[
  t = \min_{i:\, \tilde{\alpha}_i > 0} \frac{\lambda_i}{\tilde{\alpha}_i} > 0.
\]
Definiendo $\mu_i = \lambda_i - t\, \tilde{\alpha}_i \geq 0$:
\[
  x = \sum_{i=1}^{k} \mu_i v_i, \qquad \sum_{i=1}^k \mu_i = \sum \lambda_i - t\sum \tilde{\alpha}_i = 1 - 0 = 1,
\]
y por construcción de $t$, existe un índice $r$ tal que $\mu_r = 0$. Se ha reducido la representación a $k-1$ puntos. Iterando, en a lo sumo $k - (n+1)$ pasos se llega a una representación con $\leq n+1$ puntos. $\square$

\textbf{Paso 3. Conclusión.}

Combinando los dos pasos: todo $x \in P$ (poliedro acotado en $\mathbb{R}^n$) puede escribirse como combinación convexa de a lo sumo $\mathbf{n+1}$ puntos extremos de $P$. $\blacksquare$

% ============================================================
\section*{Problema 13}
\addcontentsline{toc}{section}{Problema 13}

\textbf{Enunciado.} Demuestre que el conjunto de soluciones óptimas de un programa lineal es un conjunto convexo.

\begin{marco}
\textbf{[BT §2.1] Definición de conjunto convexo:} $C \subseteq \mathbb{R}^n$ es convexo si para todo $x, y \in C$ y todo $\lambda \in [0,1]$: $\lambda x + (1-\lambda)y \in C$.

\textbf{[BT §2.1] Poliedro es convexo:} El poliedro $P = \{x \mid Ax \leq b,\, x \geq 0\}$ es convexo como intersección de semiespacios, que son conjuntos convexos.

\textbf{Linealidad de la función objetivo:} Para una función lineal $f(x) = cx$: $f(\lambda x + (1-\lambda)y) = \lambda f(x) + (1-\lambda)f(y)$.
\end{marco}

\subsection*{Solución}

Sea $Z^* = \max\{cx \mid x \in P\}$ el valor óptimo del LP (suponemos que el problema tiene solución óptima finita). Sea:
\[
  X^* = \{x \in P \mid cx = Z^*\}
\]
el conjunto de soluciones óptimas.

\textbf{Demostración de que $X^*$ es convexo.}

Sean $x, y \in X^*$ y $\lambda \in [0,1]$. Definamos $z = \lambda x + (1-\lambda)y$.

\textbf{Factibilidad de $z$:}

$x, y \in P$ y $P$ es convexo \textbf{[BT §2.1]} (como intersección de semiespacios y semiespacios de no-negatividad). Por tanto:
\[
  z = \lambda x + (1-\lambda) y \in P. \quad \checkmark
\]

\textbf{Optimalidad de $z$:}

Como $cx = cy = Z^*$ (pues $x, y \in X^*$), y usando la linealidad de $c(\cdot)$:
\[
  cz = c\bigl(\lambda x + (1-\lambda)y\bigr) = \lambda(cx) + (1-\lambda)(cy) = \lambda Z^* + (1-\lambda) Z^* = Z^*. \quad \checkmark
\]

Por tanto $z \in P$ y $cz = Z^*$, lo que implica $z \in X^*$.

Como $z = \lambda x + (1-\lambda)y \in X^*$ para todo $x, y \in X^*$ y todo $\lambda \in [0,1]$, el conjunto $X^*$ es \textbf{convexo}. $\blacksquare$

\medskip
\textbf{Observación.} El conjunto de soluciones óptimas $X^*$ es en particular un poliedro, pues es la intersección del poliedro $P$ con el hiperplano $\{x \mid cx = Z^*\}$:
\[
  X^* = \{x \mid Ax \leq b,\; x \geq 0,\; cx = Z^*\} = \{x \mid Ax \leq b,\; x \geq 0,\; cx \geq Z^*,\; -cx \geq -Z^*\},
\]
que es intersección de un número finito de semiespacios, por lo tanto un poliedro y en particular un conjunto convexo cerrado.

\end{document}
